% Páginas de entrada comunes a los libros del Preicfes Saber 11 (5 a 9)
\cleardoublepage
\thispagestyle{vacia}
{\sffamily\fontsize{26}{30}\selectfont\firalight Cómo se usa este libro\par}\vspace{2pt}
{\color{zlime}\rule{18mm}{1.6pt}}\par\vspace{8pt}

Este libro y la aplicación \marca\ son dos mitades de un mismo curso. Cada lección del libro tiene el mismo código que su
lección en la app y sigue la misma ruta: en la app la idea se construye con el dedo; aquí se mira cuadro por cuadro, se
hace con lápiz y se practica como en el examen.

\vspace{4pt}
\begin{tabularx}{\linewidth}{@{}>{\leavevmode\sffamily\small\bfseries\color{zolive}}p{34mm}X@{}}
La ruta & Abre la lección: los pasos que vas a recorrer. El último es siempre un problema.\\[4pt]
Cuadro por cuadro & Cada idea se explica en dibujos numerados, como una historieta: se mira el cuadro, se lee su texto y se
sigue.\\[4pt]
La idea & La frase que hay que llevarse de cada paso.\\[4pt]
Tu turno & Un ejercicio con el mismo dibujo. Donde la app pide arrastrar o tocar, aquí se marca con lápiz. Las soluciones
dibujadas están al final del libro.\\[4pt]
Problema resuelto & Un problema completo, resuelto paso a paso; lo que va en cada hueco está resaltado.\\[4pt]
¿Dónde está el error? & Una solución con un paso equivocado. Encontrarlo entrena a revisar tus propias cuentas.\\[4pt]
Pregunta comentada & En las pruebas de lectura, un texto con su pregunta y el análisis de por qué la respuesta es esa.\\[4pt]
Así sale en el examen & Preguntas de opción múltiple como las del Saber 11.\\[4pt]
Lo que construiste & Las ideas de la lección, para marcar las que podrías explicarle a otra persona.\\[4pt]
Evaluación & Cada mundo termina con preguntas como las del examen. Las respuestas, con su explicación, están en las claves.\\
\end{tabularx}

\vspace{8pt}
\begin{tcolorbox}[enhanced,colback=zsoft,coltext=zink,boxrule=0pt,arc=3pt,left=5mm,right=5mm,top=3mm,bottom=3mm]
{\sffamily\scriptsize\bfseries\color{zolive}UN CONSEJO}\par\vspace{3pt}
\small No pases a las preguntas sin haber hecho los «Tu turno». Una idea que entendiste mirando todavía no es tuya: se vuelve
tuya cuando la usas sola, aunque te equivoques la primera vez.
\end{tcolorbox}
\clearpage
